\section{The functional equation lemma}

In this section, we discuss the functional equation lemma for formal group laws. This lemma
can be used to construct non-isomorphic formal group laws and a universal
formal group law. \par

To begin, we need the following data. Let $K$ be a commutative ring, $R$ be a subring of $K$,
$\sigma : K \to K$ be a
ring homomorphism, $I$ be an ideal of $R$, $p$ be a prime number, $q$ be a power of this prime
number, and $s_1, s_2, \ldots $ be elements of $K$. \par
These ingredients are assumed to satisfy the following conditions:

\begin{enumerate}
    \item $\sigma (R) \subset R$,
    \item $\forall a \in R, \sigma (a) \equiv a ^ q \pmod{I}$, \label{property : frobenius}
    \item $p \in I$,
    \item $\forall i, \ s_i I \subset R$, \label{property 4}
    \item $\forall b \in K, r \in \mathbb{N}$, $I ^ r b \subset I \Rightarrow I ^ r \sigma (b) \subset I$.
        \label{condition: I_pow subset}
\end{enumerate}

\begin{definition}[{\cite[Definition 2.1.8]{hazewinkel1978formal}}]
Let $g (X) = \sum_{i = 1}^\infty b_i X ^ i $ be a power series in one variable with coefficients
in $R$. Define the \emph{recursive power series} $f_g(X)$ to be the unique solution of the following
functional equation
\begin{equation}
    f_g(X)= g(X)+ \sum_{i = 1}^\infty s_i \sigma_*^i f_g (X^{q ^ i})
    \label{equation: recurFunEq}
\end{equation}
\end{definition}

\begin{remark}
    Denote $f_g(X) = \sum _{i = 1} ^ \infty a _i X ^ i$ and let $n = q ^ r m$, where
    $q$ does not divide $m$.
    By comparing coefficients on both sides, there is a recursive formula for the
    coefficients of $f_g(X)$:
    \begin{equation}
        a _n = b _n + s_1 \sigma(a_{n/ q}) + \cdots + s _r \sigma^ r(a _{n / q ^ r}) =
        b_n + \sum_ {i = 1} ^ r s_i \sigma^i (a _ {n / q ^ i})
        \label{eq: coeff_RecurFun}
    \end{equation}

    In particular, if $q$ does not divide $n$, then $a_n = b _n$.
\end{remark}

% The functional equation lemma concerns the integrality of the coefficients of
% $f_g^{-1}(f_g(X) + f_g(Y))$.

\begin{theorem}[Functional Equation Lemma, {\cite[Chap.~1, Theorem 2.2]{hazewinkel1978formal}}]
Let $g (X)$, $h(X)$ be
two power series in one variable over $R$ and suppose that the coefficient of $X$ in $g(X)$
is invertible in $R$. Then
\begin{enumerate}
  \item the power series $F_g(X, Y) = f^{-1}_g(f_g(X) + f _g(Y))$ has its coefficients in $R$,
  \item the power series $f_g^{-1}(f_{h}(X))$ has its coefficients in $R$,
  \item there is a power series
  $\hat{h}(X)$ over $R$ such that
  $f _g (h(X)) = f _{\hat{h}}(X)$,
  \item let $\alpha(X)\in R [\![ X]\!]$, $\beta(X) \in K [\![X]\!]$, and let $r$ be a positive
    natural number, then we have
    \begin{align*}
    \alpha (X) &\equiv \beta (X) \pmod{I^r R[\![ X]\!]} \\
    \Leftrightarrow f_g(\alpha (X)) &\equiv f_g(\beta (X)) \pmod{I^r R[\![ X]\!]}.
    \end{align*}
\end{enumerate}
\label{thm: fun_eq}
\end{theorem}

\begin{remark}
    Due to equation (\ref{eq: coeff_RecurFun}), we have that the coefficient of $X$ in
    $f_g(X)$ and $g(X)$ is the same, denoted as $a_1$. Thus, the coefficient of $X$ in
    $f_g(X)$ is invertible in $R$ and
    the power series $f_g^{-1} (X)$ is well defined by Lemma \ref{lem: substInv}. By computations,
    we have that the coefficient of $X$ and $Y$ in $F_g(X,Y)$ is $a_1^{-1} a_1 = 1$, since
    the coefficient of $X$ in $f_g^{-1} (X)$ is $a_1^{-1}$.
    Moreover, we notice that
    \begin{align*}
        F_g(F_g(X,Y),Z) = f_g^{-1}(f_g(F_g(X,Y)) + f_g(Z)) &=
            f_g^{-1}(f_g(f_g^{-1} (f_g(X) + f_g(Y))) + f_g(Z))\\
             &= f_g^{-1}(f_g(X) + f_g(Y) + f_g(Z)).
    \end{align*}
    Similarly,
    \begin{align*}
        F_g(X, F_g(Y,Z)) = f_g^{-1}(f_g(X) + f_g(F_g(Y,Z))) & =
            f_g^{-1}(f_g(X) + f_g(f_g^{-1}(f_g(Y) + f_g(Z)))) \\
            &= f_g^{-1}(f_g(X) + f_g(Y) + f_g(Z)).
    \end{align*}
    Therefore, we deduce that $F_g(X,Y)$ is a commutative formal group law over $K$.
    Moreover, according to the functional equation lemma, $F_g(X, Y)$ is a commutative formal group law
    over the ring $R$.
    \label{rem: fun_eq_FG}
\end{remark}

Before discussing the proof of the functional equation integrality lemma, we first present several
technical lemmas.

\begin{lemma}[{\cite[Chap.~1, Lemma 2.4.1]{hazewinkel1978formal}}]
    Let $f _g(X) = \sum_{i = 1} ^ \infty a _n X ^ n$, and let $n = q ^ r m$ where $m$ is
    not divisible by $q$, then $a _n I^r \subset R$.
    \label{lemma : tech_lemma_one}
\end{lemma}
\begin{proof}
    By the assumption (\ref{condition: I_pow subset}) and induction,
    we have that for any given natural numbers $j, r$ and any $b \in K$,
    \begin{equation}
        I ^ r b \subset I \Rightarrow I ^ r \sigma ^ j (b) \subset I.
        \label{refinement property 5}
    \end{equation}
    \indent We proceed by induction on $r$ (the multiplicity of $q$ in $n$). \par
    Assume that the result holds for all numbers less than $r + 1$; now consider those $n$
    such that the multiplicity of $q$ in $n$ is $r + 1$.
    According to our recursive formula, the coefficient $a_n$ can be expressed as
    \[a _n = b_n + \sum_ {i = 1} ^ {r + 1} s_i \sigma^i (a _ {n / q ^ i}).\] Combining
    the result (\ref{refinement property 5}) and the assumption (\ref{property 4}), it is enough to
    prove
    \[\forall i \in [1, r + 1], \ a_{n / q ^i} I ^ {r + 1} \subset I.\]
    According to the induction hypothesis and the fact that the multiplicity of $q$
    in ${n/ q^i}$ is less than $r + 1$,
    we have $\forall i \in [1, r + 1]$, $a_{n / q ^i} I ^ {r} \subset R$. This implies that
    $a_n \in I$.
\end{proof}

\begin{lemma}
    Let $S$ be a commutative ring, $q$ be a power of a prime number $p$, for any $x,y \in S$,
    then there is an element $a \in S$ such that
    \[ (x + y) ^ q = x ^ q + y ^ q + p x y a. \]
    \label{lem: exp_prime_pow}
\end{lemma}

\begin{proof}
    Note that
    \[(x + y) ^ q = x ^ q + y ^ q + \sum_{i=1}^{q-1} \binom{q}{i} x ^ i y ^ {q-i}.  \]
    Then the lemma follows from the fact that $p \mid \binom{q}{i}$ for all $i = 1, \ldots, n-1$.
\end{proof}


\begin{lemma}
    Let $F(X_J), G(X_J) \in R[\![X_J]\!]$, where $J$ is a finite index set, let $r$ be the
    multiplicity of $q$ in $n$, namely $n = q ^ r m$ with $q \nmid m$, and let $i$ be a positive
    natural number. Assume that $F(X_J) \equiv G(X_J) \pmod{I^i}$, then we have
    \begin{equation}
        F(X_J)^n \equiv G(X_J)^n \pmod{I ^{i + r}}.
    \end{equation}
    \label{lemma: tech_lem pow_congruence}
\end{lemma}

\begin{proof}
    We prove that $F(X_J)^{q^r} \equiv G(X_J)^{q^r} \pmod{I ^{i + r}}$, from which then the
    lemma follows by taking $m$-th powers on both sides. \par
    We do induction on $r$. The case for $r =0$ is trivial. Assume that
    the result holds for $r$, namely $F(X_J)^{q^r} \equiv G(X_J)^{q^r} \pmod{I ^{i + r}}$.
    Therefore, $F(X_J) ^ {q^r} = G(X_J)^{q^r} + H(X_J)$, where all coefficients of $H(X_J)$
    are in $I^{i + r}$. Due to Lemma \ref{lem: exp_prime_pow}, there is a power series
    $\alpha (X_J) \in R [\![X_J]\!]$ such that
    \[ (G(X_J)^{q^r} + H(X_J))^q
        = (G(X_J)^{q^r})^q + H(X_J) ^ q + p G(X_J)^{q^r} H(X_J) \alpha (X_J). \]
    Then we have the following equalities
    \begin{align*}
        F(X_J)^{q^{r+1}} = (F(X_J)^{q^r})^q &= (G(X_J)^{q^r} + H(X_J))^q \\
        &= (G(X_J)^{q^r})^q + H(X_J) ^ q + p G(X_J)^{q^r} H(X_J) \alpha (X_J)
    \end{align*}
    Our result is given by the following two congruences:
    $H(X_J)^q \equiv 0 \pmod{I ^{i + r + 1}}$ and $p H(X_J) \equiv 0 \pmod{I ^{i + r + 1}}$.
\end{proof}

\begin{lemma}[Generalization of {\cite[Chap.~1 Lemma 2.4.2]{hazewinkel1978formal}}]
    Let $n,l$ be natrual numbers, $q$ be a power of prime number $p$, $J$ be a finite index set, and
    let $n=q^r m$ with $q \nmid m$. For all $G(X_J) \in R [\![X_J]\!]$, we have
    \[G (X_J)^{nq ^ l} \equiv (\sigma _*^l G (X_J^{q^l})) ^ n \pmod{I ^{r + 1}}.\]

    \label{lem: tech_lem_two}
\end{lemma}

\begin{proof}
    The case for $I = R$ is trivial. Now, we assume that $I \subsetneq R$. Then $R / I$ has
    characteristic $p$, since $p \in I$. Given property \ref{property : frobenius}, then we have
    \begin{equation*}
        G(X_J) ^{q^l} \equiv \sigma^l_*G(X_J^{q^l}) \pmod{I}
    \end{equation*}

    %%%%%%%%% Here more detail: TODO.
    We obtain the goal congruence from the above congruence and Lemma \ref{lemma: tech_lem pow_congruence}.
\end{proof}

\begin{proof}[Proof of the Functional Equation Lemma \ref{thm: fun_eq}]
For simplicity, we write $F (X, Y)$ and $f (X)$ for $F_g(X, Y)$ and $f_g(X)$ respectively. Let $F_i(X,Y)$ be the
homogeneous component of $F(X,Y)$ of degree $i$. We denote
$g(X) = \sum_{i = 1}^{\infty} b_i X ^ i$ and $f(X) = \sum_{n=1}^{\infty}a_n X ^ n$ throughout
the proof of this theorem.

\textit{Proof of part (1):}
Since $f(X) \equiv b_1 X \pmod{\deg 2}$ and $f^{-1}(X) \equiv b_1^{-1} X \pmod{\deg 2}$, we have
$F(X, Y) \equiv X + Y \pmod{\deg 2}$. This implies that the coefficients of $F_1(X, Y)$ lie in $R$.
We now proceed by induction to prove that $F_n(X, Y) \in R[\![X, Y]\!]$.

Suppose that $F_1(X, Y), \dots, F_{n - 1}(X, Y) \in R[\![X, Y]\!]$. Note that for $r \ge 2$, we have
\[ (F(X, Y))^r \equiv (F_1(X, Y) + \cdots + F_{n- 1}(X, Y))^r \pmod{\deg(n + 1)}. \]
This follows from the fact that the difference between the expressions inside the
parentheses consists only of terms of degree at least $n$. Upon taking the $r$-th power, any
term in the expansion involving this difference will have a degree of at least $n+1$.

By the definition of $F(X, Y) = f^{- 1} (f (X) + f (Y))$, we have
\begin{align*}
    f (F (X, Y)) &= f (X) + f (Y)  \\
    \sigma_*^i f(\sigma_*^i F(X, Y)) &= \sigma_*^i f (X) + \sigma_*^i f (Y).
\end{align*}
Recall that $f$ and $g$ satisfy recursion formula (\ref{equation: recurFunEq}),
\[f (X) = g (X) + \sum_ {n = 1} ^ \infty s_n\sigma _*^n f (X ^ {q ^ n}).\]
By substituting $F(X, Y)$ into $X$ and writing $f(X)=\sum _{n = 1} ^ \infty a_n X ^ n$,
we obtain that
\begin{align*}
    f (F(X, Y)) & = g(F(X,Y)) + \sum_{n=1}^{\infty}s_n \sigma^n_* f(F(X,Y)^{q^n}) \\
            & = g (F(X, Y)) + \sum _{i = 1} ^ \infty s_i \sum _{n = 1} ^ \infty
            \sigma ^ i (a _n) (F(X, Y))^{q ^i n}.
\end{align*}

Next, by Lemma \ref{lem: tech_lem_two}, we have
\[s_i \sigma^i(a_n) F(X, Y)^{q^i n} \equiv s_i\sigma^i(a_n) (\sigma_*^i F(X^{q^i}, Y^{q^i}))^n
\pmod{(R, \deg(n + 1))}. \]

The notation $\pmod{(R, \deg(n + 1))}$ indicates that the two expressions are congruent up to terms of degree $n$, and that their corresponding coefficients are congruent modulo the subring $R$. Then, modulo $(R, \deg(n + 1))$, we have
\[ \begin{aligned}
    f(F(X, Y)) & \equiv g(F(X, Y)) + \sum_{i = 1}^\infty s_i \sum_{n = 1}^\infty \sigma^i(a_n) (\sigma_*^i F(X^{q^i}, Y^{q^i}))^n \\
    & \equiv g(F(X, Y)) + \sum_{i = 1}^\infty s_i \sigma_*^i f (\sigma_*^i F(X^{q^i}, Y^{q^i})) \\
    & \equiv g(F(X, Y)) + \sum_{i = 1}^\infty s_i (\sigma_*^i f(X^{q^i}) + \sigma_*^i f(Y^{q^i})) \\
    & \equiv g(F(X, Y)) + f(X) + f(Y) - g(X) - g(Y).
\end{aligned} \]

Since $g(X) \equiv b_1 X \pmod{\deg 2}$, then
\[g (F(X, Y)) \equiv b_1F_n (X, Y) \pmod{(R, \deg (n + 1))}. \]

Using the equation $f(F(X,Y)) = f (X) + f (Y)$ and the congruence above, we have
\begin{align*}
    f(X) + f(Y) &= f(F(X,Y)) \\
        &\equiv g(F(X,Y)) + f(X) + f(Y) - g(X) - g(Y)\\
        &\equiv b_1 F_n(X,Y) + f(X) + f(Y) - g(X) - g(Y) \pmod{(R, \deg (n + 1))}.
\end{align*}

Since $g(X) \in R[\![X]\!]$, we have $g(X) \equiv 0 \pmod{R}$. This implies that
\[b_1 F_n(X, Y) \equiv 0 \pmod{(R, \deg(n + 1))}. \]

Because $b_1$ is invertible in $R$, we conclude that $F_n(X, Y) \equiv 0 \pmod{(R, \deg(n + 1))}$.
By induction, the homogeneous component of degree $n$ lies in $R[\![X, Y]\!]$ for all $n$.
It follows that $F(X, Y) \in R[\![X, Y]\!]$.

\textit{Proof of part (2):} Write $G(X) = f_g^{-1}(f_{h}(X))$. Let $G_i(X)$ be the homogeneous component of degree $i$ in $G(X)$. We proceed by induction to prove that for all $i$, the coefficients of $G_i(X)$ lie in $R$. Assume that this result holds for all $i < n$. Now consider the case $i = n$. Modulo $(R, \deg(n + 1))$, we have
\[ \begin{aligned}
    f_{h}(X) &= f(G(X)) \\
    &= g(G(X)) + \sum_{i = 1}^\infty s_i \sigma_*^i f(G(X)^{q^i}) \\
    &= g(G(X)) + \sum_{i = 1}^\infty s_i \sigma_*^i \sum_{m = 1}^\infty a_m G(X)^{q^i m} \\
    &\equiv g(G(X)) + \sum_{i = 1}^\infty s_i \sum_{m = 1}^\infty \sigma_*^i(a_m) (\sigma_*^i G(X^{q^i}))^m \\
    &= g(G(X)) + \sum_{i = 1}^\infty s_i \sigma_*^i f(G(X^{q^i})) \\
    &= g(G(X)) + \sum_{i = 1}^\infty s_i \sigma_*^i (f_{h}(X^{q^i})) \\
    &= g(G(X)) + f_{h}(X) - h(X).
\end{aligned} \]

It follows that $g(G(X)) \equiv h(X) \pmod{(R, \deg(n + 1))}$. We then have
\[h(X) \equiv g(G(X)) \equiv b_1 G_n(X) \pmod{(R, \deg(n + 1))}.\]
Because $b_1$ is invertible, we conclude that the coefficients of $G_n(X)$ lie in $R$.

\textit{Proof of part (3):} Let $\hat{f}(X) = f(h(X))$. By the functional equation \eqref{equation: recurFunEq}, it suffices to show that the coefficients of $\hat{f}(X) - \sum_{i=1}^\infty s_i \sigma_*^i \hat{f}(X^{q^i})$ lie in $R$. Working modulo $R$, we have the following sequence of equalities and congruences:
\[\begin{aligned}
    \hat{f}(X) - \sum_{i=1}^\infty s_i \sigma_*^i\hat{f}(X^{q^i}) &= f(h(X)) - \sum_{i=1}^\infty s_i \sigma_*^i f(\sigma_*^i h(X^{q^i})) \\
    &= f(h(X)) - \sum_{i=1}^\infty s_i \sum_{n=1}^\infty \sigma^i(a_n)(\sigma_*^i h(X^{q^i}))^n \\
    &\equiv f(h(X)) - \sum_{i=1}^\infty s_i \sum_{n=1}^\infty \sigma^i(a_n)(h(X)^{q^i n}) \\
    &= f(h(X)) - \sum_{i=1}^\infty s_i \sigma_*^i f(h(X)^{q^i}) \\
    &= g(h(X)) \equiv 0.
\end{aligned} \]
Since the expression reduces to $0$ modulo $R$, we conclude that all its coefficients are indeed in $R$.

\textit{Proof of part (4):}
For the forward implication, assume that $\alpha(X) \equiv \beta(X) \pmod{I^r R[\![X]\!]}$.
Then $\beta(X) \in R[\![X]\!]$. By Lemma \ref{lemma: tech_lem pow_congruence},
we have $\alpha(X)^n \equiv \beta(X)^n \pmod{I^{r + t}}$, where $t$ is the multiplicity of
$q$ in $n$. Denote $f(X) = \sum_{i=1}^\infty a_i X^i$. By Lemma \ref{lemma : tech_lemma_one},
we have $a_n\alpha(X)^t \equiv a_n\beta(X)^t \pmod{I^r}$, where $t$ is the multiplicity of
$q$ in $n$. Therefore, $f(\alpha(X)) \equiv f(\beta(X)) \pmod{I^r}$.

For the reverse implication, assume that $f(\alpha(X)) \equiv f(\beta(X)) \pmod{I^r R[\![X]\!]}$.
First, we prove the following result:
\begin{equation}
    \alpha(X) \equiv 0 \pmod{I^r} \implies f^{-1}(\alpha(X)) \equiv 0 \pmod{I^r}.
    \label{eq: implication_f_inv}
\end{equation}

Let $\gamma(X) = f^{-1}(\alpha(X))$; then $\alpha(X) = f(\gamma(X))$. Let $\gamma_n(X)$ be
the truncation of $\gamma(X)$ at degree $n + 1$. We proceed by induction to prove that
$\gamma_n(X) \in I^r R[\![X]\!]$ for all $n$. The cases for $k = 1, 2$ are trivial.
Assume that the result holds for $k$, namely $\gamma_k(X) \equiv 0 \pmod{I^r}$.
Now consider $\gamma_{k+1}(X)$. By the functional equation and the definition of $\gamma(X)$,
we have

\begin{equation}
    \alpha(X) = g(\gamma(X)) + \sum_{i=1}^\infty s_i \sigma_*^i f(\gamma(X)^{q^i}).
    \label{eq: decomp alpha gamma}
\end{equation}

We claim that:
\begin{equation}
    \sum_{i=1}^\infty s_i \sigma_*^i f(\gamma(X)^{q^i}) \in I^r R[\![X]\!] \pmod{\deg(k + 1)}.
    \label{eq: claim_trunc}
\end{equation}

Comparing the coefficients of degree $k+1$ on both sides of equation \eqref{eq: decomp alpha gamma},
we obtain $b_1 c_{k+1} \in I^r$, where $\gamma(X) = \sum_{i=1}^\infty c_i X^i$ and
$g(X) = \sum_{i=1}^\infty b_i X^i$. Since $b_1$ is a unit in $R$, it follows that $c_{k+1} \in I^r$.

\textit{Proof of Claim \eqref{eq: claim_trunc}:} By Lemma \ref{lemma: tech_lem pow_congruence}
and the induction hypothesis, we have $\gamma_k(X)^n \equiv 0 \pmod{I^{r + t}}$,
where $t$ is the multiplicity of $q$ in $n$. By assumption \eqref{property 4},
it suffices to show that

\[\sigma_*^i f(\gamma(X)^{q^i}) \in I^{r+1} R[\![X]\!] \pmod{\deg(k + 2)},\]

which is equivalent to

\begin{equation}
    \sigma^i(a_n) \gamma(X)^{q^i n} \in I^{r+1} R[\![X]\!] \pmod{\deg(k + 2)}.
    \label{eq: sigma_gamma_mem}
\end{equation}

Since $\gamma(X)^{q^i n} \equiv \gamma_k(X)^{q^i n} \pmod{\deg(k + 2)}$, it follows from
Lemma \ref{lemma : tech_lemma_one} that the congruence \eqref{eq: sigma_gamma_mem} holds.
This proves the implication \eqref{eq: implication_f_inv}.

By assumption, $f(\beta(X)) - f(\alpha(X)) \in I^r R[\![X]\!]$. Thus, applying
\eqref{eq: implication_f_inv}, we have
\[\delta(X) := f^{-1}(f(\beta(X)) - f(\alpha(X))) \in I^r R[\![X]\!]. \]
Then
\[\begin{aligned}
    f(\beta(X)) &= f(\delta(X)) + f(\alpha(X)), \quad \text{so that} \\
    \beta(X) &= f^{-1}(f(\delta(X)) + f(\alpha(X))).
\end{aligned} \]

By Remark \ref{rem: fun_eq_FG}, $f^{-1}(f(X) + f(Y))$ is a formal group law. Hence, we can write

\[\beta(X) = f^{-1}(f(\delta(X)) + f(\alpha(X))) = \alpha(X) + \delta(X) G(\delta(X), \alpha(X))\]

for some $G(X, Y) \in R[\![X, Y]\!]$. Since $\delta(X) \in I^r R[\![X]\!]$.
We conclude that
\[\beta(X) \equiv \alpha(X) \pmod{I^r R[\![X]\!]}.\]



% \textit{Proof of part 4:} For the forward implication: assume that
% $\alpha (X) \equiv \beta (X) \pmod{I^r R[\![ X]\!]}$, then $\beta (X) \in R [\![X]\!]$.
% According to the lemma \ref{lemma: tech_lem pow_congruence}, we have that
% $\alpha (X)^n\equiv \beta (X)^n \pmod{I^{r + t}}$, where $t$ is the multiplicity of $q$ in $n$.
% Let $f(X) = \sum_{i=1}^\infty a_i X ^i $,  according to the lemma \ref{lemma : tech_lemma_one},
% we have that $a_n\alpha(X)^t \equiv a_n\beta(X)^t \pmod{I ^ r}$, where $t$ is multiplicity
% of $q$ in $n$. Therefore, $f(\alpha (X)) \equiv f(\beta(X)) \pmod{I ^ r}$.

% For the reverse direction, assume that $f(\alpha (X)) \equiv f(\beta (X)) \pmod{I^r R[\![ X]\!]}$.

% First, we prove the following result:

% \begin{equation}
%     \alpha (X) \equiv 0 \pmod{I ^ r} \Rightarrow f^{-1}(\alpha (X)) \equiv 0 \pmod{I ^ r}
%     \label{eq: implication_f_inv}
% \end{equation}

% Let $\gamma(X) = f^{-1}(\alpha (X))$, then $\alpha (X) = f(\gamma (X))$. Let $\gamma_n (X)$ be
% the truncation of degree $n + 1$ of $\gamma (X)$, and we use induction to prove all
% $\gamma_n(X) \in I^rR[\![X]\!]$. The cases $k = 1, 2$ are trivial. Assume that the result
% holds for $k$, namely $\gamma_k(X) \equiv 0 \pmod{I ^ r}$; now we consider $\gamma_{k+1}(X)$.
% According to the functional equation and definition of $\gamma (X)$, we have

% \begin{equation}
%     \alpha(X) = g(\gamma(X)) + \sum_{i=1}^\infty s_i \sigma ^ i_* f(\gamma(X) ^ {q^i})
%     \label{eq: decomp alpha gamma}
% \end{equation}

% We claim that:

% \begin{equation}
%     \sum_{i=1}^\infty s_i \sigma ^ i_* f(\gamma(X) ^ {q^i}) \in I ^ r R[\![X]\!] \pmod{\deg (k + 1)}.
%     \label{eq: claim_trunc}
% \end{equation}

% Considering the $(k+1)$-th coefficient on both sides of equation \ref{eq: decomp alpha gamma}, we
% have $b_1 * c_{k+1} \in I ^ r$, where $\gamma(X) = \sum_{i=1}^\infty c_i X^i$ and
% $g(X) = \sum_{i=1}^\infty b_i X ^ i$. Since $b_1$ is a unit in $R$, then $c_{k+1} \in I ^ r$.

% Proof of claim \ref{eq: claim_trunc} above. According to the
% lemma \ref{lemma: tech_lem pow_congruence} and the hypothesis, we have
% $\gamma_k(X) ^n \equiv 0 \pmod{I ^{r + t}}$, where $t$ is multiplicity of $q$ in $n$.
% By the assumption \ref{property 4}, it is enough to prove that

% \begin{equation}
%     \sigma^i_*f(\gamma(X)^{q^i}) \in I ^ {r+1} R[\![X]\!] \pmod{\deg (k + 2)}
% \end{equation}

% which is equivalent to

% \begin{equation}
%     \sigma^i (a_n) \gamma(X)^{q^i n} \in I^{r+1} R [\![X]\!] \pmod{\deg (k + 2)}
%     \label{eq: sigma_gamma_mem}
% \end{equation}

% Since
% \[\gamma (X)^{q^in} \equiv \gamma_k (X) ^ {q^i n} \pmod{\deg (k + 2)}\] and
% lemma \ref{lemma : tech_lemma_one}, we have the congruence \ref{eq: sigma_gamma_mem}.
% This proves the result \ref{eq: implication_f_inv}.

% By assumption, $f(\alpha (X)) - f(\beta(X)) \in I^r R[\![X]\!]$, so we have
% \[\delta (X) := f^{-1} (f(\beta (X)) - f(\alpha(X))) \in I^r R[\![X]\!].\]
% Then

% \begin{align*}
%     f (\beta (X)) &=  f(\delta (X)) + f(\alpha (X))\\
%     \beta (X) & = f^{-1} (f(\delta (X)) + f(\alpha (X)))
% \end{align*}

% It is straightforward to prove that $f^{-1}(f(X) + f(Y))$ is a formal group law. Hence,

% \[\beta (X) = f^{-1} (f(\delta (X)) + f(\alpha (X))) = \alpha (X) + \delta (X) G(\delta(X), \alpha (X)).\]

% for some $G(X,Y) \in R [\![X,Y]\!]$. Since $\delta(X) \in I^r R[\![X]\!]$,
% $\beta (X) \equiv \alpha (X) \pmod{I^r R[\![X]\!]}$.


\end{proof}